\documentclass[conference,a4paper]{IEEEtran}

\usepackage[turkish]{babel}
\usepackage[utf8]{inputenc}
\usepackage[T1]{fontenc}
\usepackage{graphicx}
\usepackage{float}
\usepackage{hyperref}
\usepackage{amsmath}
\usepackage{microtype} 
\usepackage{booktabs}  
\usepackage{xcolor}    
\usepackage{tikz}      
\usetikzlibrary{shapes.geometric, arrows.meta}
\usepackage[linesnumbered,ruled,vlined]{algorithm2e} 

\SetKwInput{KwData}{Girdi}
\SetKwInput{KwResult}{Çıktı}
\SetAlgorithmName{Algoritma}{algoritma}{Algoritmalar Listesi}

\addto\captionsturkish{
  \renewcommand{\abstractname}{\textbf{Özet}}
  \renewcommand{\figurename}{Şekil}
  \renewcommand{\tablename}{Tablo}
  \renewcommand{\refname}{Kaynaklar}
}

\hypersetup{
    colorlinks=true,
    linkcolor=blue,
    filecolor=magenta,      
    urlcolor=blue,
    pdftitle={Kelime Bulmaca Proje Raporu},
}

\begin{document}

% Turkish Babel paketinin TikZ ve Algorithm paketleriyle çakışmasını engellemek için:
\shorthandoff{=}

\title{Yazılım Laboratuvarı II - Proje 2\\Kelime Bulmaca (Word Crush) Oyunu Geliştirme Raporu}

\author{\IEEEauthorblockN{Öğrenci Adı Soyadı}
\IEEEauthorblockA{\textit{Bilgisayar Mühendisliği Bölümü} \\
\textit{Kocaeli Üniversitesi}\\
Öğrenci No: XXXXXXXXX \\
email@ogrenci.edu.tr}
}

\maketitle

\begin{abstract}
Bu raporda, Yazılım Laboratuvarı II dersi Proje 2 kapsamında tasarlanan ve geliştirilen Kelime Bulmaca (Word Crush) oyununun tüm mühendislik aşamaları sunulmaktadır. Proje, başlangıçta planlanan geleneksel veritabanı bağımlı yapısından arındırılarak modern mobil oyunların "Offline-First" (Öncelikle Çevrimdışı) yaklaşımına uygun bir mimariye dönüştürülmüştür. Sunucu tarafında C\# ve ASP.NET Core, verileri tutan değil; yüksek işlem gücü gerektiren kelime doğrulama, Trie ağacı aramaları ve yerçekimi algoritmalarını işleyen bağımsız bir "Algoritma Motoru (Stateless Engine)" olarak yapılandırılmıştır. İstemci tarafında React Native kullanılmış olup, kullanıcı verileri ve skor tabloları AsyncStorage üzerinden cihazın yerel hafızasına (Local Storage) entegre edilmiştir. Rapor boyunca; donanım ivmeli partikül patlama efektleri (Particle System), Glassmorphism arayüz tasarımları, logaritmik arama hızına sahip veri yapıları ve yerel depolama optimizasyonları detaylıca incelenmektedir.
\end{abstract}

\begin{IEEEkeywords}
Mobil Oyun, React Native, C\#, AsyncStorage, Trie Veri Yapısı, Particle System, Offline-First Mimari.
\end{IEEEkeywords}

\section{Giriş}
Mobil platformlar, günümüzde eğlence ve zihinsel gelişim araçlarının en önemli taşıyıcısı konumundadır. Akıllı cihazların donanımsal kapasitelerinin artması, karmaşık algoritmalar gerektiren oyunların dahi mobil cihazlarda akıcı bir şekilde çalıştırılabilmesine olanak tanımaktadır. 

Bu projenin temel hedefi, sıradan bir kelime eşleştirme oyununun ötesine geçerek, içerisinde gelişmiş kombo mekanikleri, zorluk seviyelerine göre dinamik olarak üretilen rastgele harf matrisleri ve oyun içi satın alımlara dayalı yetenek (joker) sistemleri barındıran modern bir "Word Crush" oyunu geliştirmektir.

Geliştirme sürecinin ilerleyen safhalarında, oyunun her an ve her yerde internet bağlantısına ve dış bir veritabanı sunucusuna muhtaç olmadan oynanabilmesi gerekliliği ortaya çıkmıştır. Bu doğrultuda mimari büyük bir evrim geçirmiş; sunucu sadece ağır algoritmik iş yükünü (Backend as an Algorithm Engine) omuzlarken, durum yönetimi (State Management) ve veri kalıcılığı (Data Persistence) tamamen React Native uygulaması üzerine (AsyncStorage ile) kaydırılmıştır.

\section{Sistem ve Yazılım Mimarisi}
Geliştirilen sistem, bağımsız modüller halinde tasarlanmış ve dağıtık veri sorumluluğu (Distributed Data Responsibility) prensibine göre inşa edilmiştir.

\subsection{Algoritma Motoru ve Yerel Depolama (Stateless - Offline First) Modeli}
Oyunun genel mimarisi Şekil \ref{fig:mimari}'de gösterildiği üzere iki ana bileşenden oluşmaktadır: Yerel depolamaya sahip Mobil İstemci (React Native) ve Durumsuz Sunucu Uygulaması (Stateless C\# API).

\begin{figure}[H]
\centering
\begin{tikzpicture}
% Absolute positioning kullanarak hataları önledik
\node[rectangle, draw=blue, fill=cyan, very thick, minimum width=6.5cm, minimum height=1.2cm, text centered, rounded corners=5pt, font=\sffamily] (app) at (0, 3) {\textbf{İstemci (React Native) + AsyncStorage}};

\node[rectangle, draw=orange, fill=yellow, very thick, minimum width=6.5cm, minimum height=1.2cm, text centered, rounded corners=5pt, font=\sffamily] (api) at (0, 0.5) {\textbf{Algoritma Motoru (Stateless C\# API)}};

% Oklar
\draw [<->, thick] (app) -- node[right, font=\footnotesize] {Oyun İçi Hamleler (REST / JSON)} (api);
\end{tikzpicture}
\caption{Yerel Depolama ve Durumsuz Motor Mimarisi}
\label{fig:mimari}
\end{figure}

Bu yapının en büyük avantajı, sunucunun veritabanı bağlantı havuzları (Connection Pooling) ile yorulmaması ve istemcinin çevrimdışı veriye anında erişebilmesidir.

\subsection{Arka Uç (Backend) Mimari Katmanları}
C\# ile geliştirilen arka uç uygulaması, veritabanı katmanı (Entity Framework/Repository) tamamen silinerek, sadece matematiksel ve dizgisel (string) hesaplamalara odaklanmış bir motora dönüştürülmüştür:
\begin{itemize}
    \item \textbf{Sunum Katmanı (Controllers):} İstemciden gelen hamle koordinatlarını ve seçilen kelimeleri kabul eden uç noktalardır.
    \item \textbf{İş Mantığı Katmanı (Services):} Uygulamanın kalbidir. Kelimelerin Trie ağacı üzerinden doğrulandığı, boşalan hücrelerin yerçekimi (Gravity) ile doldurulduğu ve matris kilitlenmelerinin (Deadlock) çözüldüğü bölümdür.
    \item \textbf{Önbellek (In-Memory Cache):} İstemci kopukluklarına karşı, oyunun anlık matris dizilimi sunucunun geçici RAM hafızasında \texttt{IMemoryCache} aracılığıyla güvende tutulur.
\end{itemize}

\subsection{Kullanılan Teknolojiler}
Tablo \ref{tab:teknolojiler}'de sistemin inşasında kullanılan güncel teknolojiler özetlenmiştir.

\begin{table}[H]
\centering
\caption{Projeye Entegre Edilen Teknolojiler ve Görevleri}
\label{tab:teknolojiler}
\resizebox{\columnwidth}{!}{%
\begin{tabular}{@{}lll@{}}
\toprule
\textbf{Sistem Bölümü} & \textbf{Kullanılan Teknoloji} & \textbf{Temel Kullanım Amacı} \\ \midrule
\textbf{Ön Uç (İstemci)} & React Native, Animated API & Arayüz ve Görsel Efektler \\
\textbf{Yerel Depolama}  & AsyncStorage & Veritabanı Bağımsız Kalıcılık \\
\textbf{Arka Uç (Sunucu)}& C\#, ASP.NET Core & Algoritma ve Puanlama Motoru \\
\textbf{Algoritmik Altyapı}& Trie Ağacı, Fisher-Yates & O(L) Doğrulama, Karma \\
\textbf{Görsel Şekillendirme}& Particle System, Glassmorphism & Zengin Etkileşim Tasarımı \\ \bottomrule
\end{tabular}%
}
\end{table}

\section{Kullanıcı Arayüzü ve UX Tasarımı}
Kullanıcı deneyimi (UX), modern donanımların grafik işleme gücünden (GPU) tam faydalanacak şekilde tasarlanmıştır. React Native \texttt{Animated} API ve \texttt{LayoutAnimation} özellikleri sıkça kullanılmıştır.

\subsection{Gelişmiş Görsel Efektler ve Animasyonlar}
Oyun içi etkileşimi maksimize etmek amacıyla standart tasarımların ötesine geçilmiştir:
\begin{itemize}
    \item \textbf{Parçacık Sistemi (Particle System):} Oyuncu geçerli bir kelimeyi eşleştirdiğinde veya özel bir yetenek (Joker) kullandığında, o noktanın merkezinden etrafa kavisli ve parabolik olarak saçılan parçacık (yıldız ve küre) efektleri tetiklenir.
    \item \textbf{Yüzen Metinler (Floating Text):} Kelimelerin patlamasıyla eşzamanlı olarak kazanılan puan veya kombo mesajı (Örn: \texttt{+150}) ekranda yukarı doğru süzülerek kaybolur.
    \item \textbf{Seçim Hissi (Scale Glow):} Parmağın sürüklendiği harfler statik kalmaz; cihaz donanım ivmelendirmesiyle (Native Driver) büyür (scale) ve etrafında altın rengi bir parlama (glow) oluşturur.
\end{itemize}

\begin{figure}[H]
    \centering
    \includegraphics[width=0.85\linewidth]{oyun_ekrani.png}
    \caption{Dinamik Harf Matrisi, Skor Tablosu ve Özel Joker Yetenekleri}
    \label{fig:oyun_ekrani}
\end{figure}

\section{Algoritmik Yaklaşımlar ve Veri Yapıları}
Sistemin donmadan çalışabilmesi için arka uçta çalışan kelime doğrulama algoritmalarının minimum karmaşıklıkta işlemesi gereklidir.

\subsection{Trie (Önek Ağacı) ile Kelime Doğrulama}
Oyuncunun tahtadan seçtiği bir karakter dizisinin geçerli bir kelime olup olmadığını kontrol etmek için \texttt{Trie (Prefix Tree)} veri yapısı implemente edilmiştir. Klasik lineer aramanın aksine, Trie yapısı arama hızını sadece aranan kelimenin uzunluğu olan $O(L)$ karmaşıklığına düşürür.

\subsection{Gelişmiş Kombo ve Skor Hesaplama}
Oyun, seçilen ana kelimenin içindeki gizli alt kelimeleri de (sub-words) bulup puanlandırır. Algoritma \ref{alg:kombo}, ana kelime içerisindeki tüm alt harf dizilimlerini üreterek geçerli olanlara ekstra kombo puanı ekleyen yapıyı özetlemektedir.

\begin{algorithm}[H]
\caption{Trie Üzerinde Kelime ve Kombo Doğrulama}
\label{alg:kombo}
\KwData{secilenKelime (String), trieKoku (TrieNode)}
\KwResult{toplamPuan, bulunanKelimeler Listesi}
$toplamPuan \gets 0$\;
$bulunanKelimeler \gets \text{Bos Liste}$\;

\tcp{Önce ana kelime kontrol edilir}
\If{\text{TrieAra}(secilenKelime, trieKoku)}{
    $toplamPuan \gets \text{UzunlugaGoreHesapla}(secilenKelime)$\;
    $bulunanKelimeler.\text{Ekle}(secilenKelime)$\;
}

\tcp{Ardından kelime içerisindeki alt kombinasyonlar aranır}
\ForEach{altKelime $\in$ \text{AltKelimeleriUret}(secilenKelime)}{
    \If{\text{Uzunluk}(altKelime) $\ge 3$}{
        \If{\text{TrieAra}(altKelime, trieKoku)}{
            \If{altKelime $\notin$ bulunanKelimeler}{
                $toplamPuan \gets toplamPuan + \text{EkstraKombo}(altKelime)$\;
                $bulunanKelimeler.\text{Ekle}(altKelime)$\;
            }
        }
    }
}
\Return{toplamPuan, bulunanKelimeler}\;
\end{algorithm}

\section{Veri Kalıcılığı ve Performans Optimizasyonları}

\subsection{Yerel Hafıza (AsyncStorage) Mimarisi}
Projeyi geleneksel web uygulamalarından ayıran en büyük özellik veritabanı bağımsızlığıdır. Kullanıcıların profil verileri, sahip oldukları toplam altınlar ve kariyer skor geçmişleri, sunucu yerine React Native'in sunduğu \texttt{AsyncStorage} API'si kullanılarak cihazın kalıcı belleğine asenkron olarak yazılır. 

Sunucu tarafı yalnızca gönderilen kelimenin matematiksel karşılığını (kaç puan getireceğini, yeni tahta dizilimini) hesaplar. İstemci bu veriyi alır ve kendi hafızasındaki (Örn: 10.000 altın) değere yansıtarak günceller.

\subsection{Performans ve Render Optimizasyonları}
100 hücreli bir matriste kaydırma yaparken FPS (Frame per Second) düşüşü yaşamamak adına hücreler \texttt{React.memo} ile sarmalanmış, yalnızca durumu (\texttt{isSelected}) değişen hücrelerin ekrana yeniden çizdirilmesi sağlanmıştır. Sunucu tarafında ise API istekleri \texttt{async/await} yapılarıyla thread-safe olarak asenkronize çalıştırılır.

\section{Sonuç ve Gelecek Geliştirmeler}
Yazılım Laboratuvarı II kapsamında tasarlanan Kelime Bulmaca oyunu; başlangıçtaki geleneksel mimarisini kırarak, modern mobil dünyasının "veritabanı bağımsız (standalone), yüksek grafik etkileşimli ve durumsuz arka uç (stateless backend)" trendlerini başarıyla uygulamıştır. Trie veri yapıları, DFS arama mekanizmaları ve gelişmiş partikül (particle) algoritmaları projenin mühendislik değerini en üst düzeye çıkarmıştır.

Gelecek geliştirmelerde, mevcut yerel depolama altyapısına opsiyonel bir "Bulut Senkronizasyonu (Cloud Sync)" eklenerek, oyuncuların farklı cihazlarda kaldıkları yerden devam etmelerine olanak sağlanması planlanmaktadır.

\vspace{1cm}

\begin{thebibliography}{10}

\bibitem{reactnative_docs}
Meta Platforms, Inc., ``React Native Documentation: Building Native Apps with React,'' [Çevrimiçi]. Erişim adresi: \url{https://reactnative.dev/docs/getting-started}.

\bibitem{aspnet_core}
Microsoft, ``ASP.NET Core Architecture and Web APIs,'' [Çevrimiçi]. Erişim adresi: \url{https://learn.microsoft.com/en-us/aspnet/core/}.

\bibitem{cormen}
T. H. Cormen, C. E. Leiserson, R. L. Rivest, ve C. Stein, \textit{Introduction to Algorithms}, 3. Baskı, Cambridge, MA, ABD: MIT Press, 2009.

\bibitem{trie_ds}
E. Fredkin, ``Trie memory,'' \textit{Communications of the ACM}, vol. 3, no. 9, sf. 490--499, 1960.

\bibitem{design_patterns}
E. Gamma, R. Helm, R. Johnson, ve J. Vlissides, \textit{Design Patterns: Elements of Reusable Object-Oriented Software}, Reading, MA, ABD: Addison-Wesley, 1994.

\bibitem{local_storage}
MDN Web Docs, ``Web Storage API / AsyncStorage Concepts,'' [Çevrimiçi]. Erişim adresi: \url{https://developer.mozilla.org/en-US/docs/Web/API/Web_Storage_API}.

\end{thebibliography}

\end{document}
